% getting_started.tex — step-by-step guide from installation to first use

\chapter{Getting Started}

This chapter guides a brand-new user from installation through the first
successful synchronisation of a nested Git repository tree.

% -----------------------------------------------------------------------
\section{Prerequisites}

\begin{itemize}
  \item Python~3.11 or later.
  \item Git~2.30 or later available on \texttt{PATH}.
  \item (Optional) PyYAML~6.x for \texttt{.yml} document support.
\end{itemize}

% -----------------------------------------------------------------------
\section{Installation}

Install \cgspkg{} from PyPI:

\begin{lstlisting}[language=bash]
$ uv pip install ComplexGitSync
\end{lstlisting}

To include optional YAML support:

\begin{lstlisting}[language=bash]
$ uv pip install "ComplexGitSync[yaml]"
\end{lstlisting}

To install from source with developer extras (needed to run the test suite):

\begin{lstlisting}[language=bash]
$ git clone https://github.com/flipoyo/ComplexGitSync.git
$ cd ComplexGitSync
$ uv sync --extra dev
\end{lstlisting}

Verify the installation:

\begin{lstlisting}[language=bash]
$ cgitsync --help
\end{lstlisting}

% -----------------------------------------------------------------------
\section{Creating a Project Spec (\texttt{.cgs})}

A \emph{local authoring spec} (\texttt{.cgs}) is a TOML file that describes
the repository tree you want to manage. For the self-standing
\texttt{ComplexGitSync} setup, create \texttt{complexgitsync.cgs}:

\begin{lstlisting}[language=toml]
[document]
format_version = "1.0"

[project]
name           = "ComplexGitSync"
default_branch = "main"

[[repos]]
gitprovider        = "github"
project_owner_name = "flipoyo"
project_name       = "ComplexGitSync"
relative_path      = "."

[[repos]]
gitprovider        = "github"
project_owner_name = "flipoyo"
project_name       = "DevSpec"
relative_path      = "DevSpec"
nested_config      = "disabled"

[[repos]]
gitprovider        = "github"
project_owner_name = "flipoyo"
project_name       = "DocComplexGitSync"
relative_path      = "docs"
nested_config      = "doccomplexgitsync.cgs"
\end{lstlisting}

In this topology, \texttt{docs/doccomplexgitsync.cgs} then declares
\texttt{DocSpec} (resolved as \texttt{docs/DocSpec}
from the repository root) as a nested documentation dependency.

% -----------------------------------------------------------------------
\section{Validating the Spec}

Check that \cgscli{} can parse and resolve the spec without errors:

\begin{lstlisting}[language=bash]
$ cgitsync validate complexgitsync.cgs
\end{lstlisting}

A successful run prints the registry summary with \texttt{lifecycle\_state:
DECLARED} and exits with code~0.

% -----------------------------------------------------------------------
\section{Inspecting the Dependency Tree}

Render the dependency graph as text:

\begin{lstlisting}[language=bash]
$ cgitsync tree complexgitsync.cgs
\end{lstlisting}

% -----------------------------------------------------------------------
\section{Next Steps}

Once the spec is validated you can proceed to:
\begin{itemize}
  \item \textbf{clone} — check out the full tree from scratch.
  \item \textbf{describe} — inspect a generated \texttt{.gts} snapshot.
\end{itemize}

Refer to Chapter~\ref{chap:user-guide} for complete command and API
documentation.
